\pagestyle{fancy}
\fancyhead[l]{\autorUO}
\fancyfoot[l]{\asignaturaAbbr}
\fancyfoot[r]{\fecha}

\section{Introducción} \label{sec:1}
En el marco de este Trabajo Fin de Máster (TFM), surge la necesidad de desarrollar soluciones computacionales de alto rendimiento aplicadas
al procesamiento digital de imágenes. Con el crecimiento exponencial de los datos y la complejidad de los algoritmos modernos, es de gran
importancia aprovechar las capacidades de las unidades de procesamiento gráfico (GPU) para acelerar las tareas computacionales intensivas.
Este documento tiene como objetivo introducir los fundamentos necesarios para comprender el uso de \pytorch\, una de las bibliotecas de
computación científica más utilizadas en la actulidad.

El enfoque principal de este documento se centrará en los tensores de \pytorch, los cuales constituyen la estructura de datos fundamental para
representar información en esta biblioteca. Entender su creación, manipulación y operaciones básicas es imprescindible para avanzar hacia
algoritmos más complejos.








